\documentclass[UTF8,a4paper,11pt]{ctexart}

% ── Page geometry ──
\usepackage[margin=2.2cm,top=2.5cm,bottom=2.5cm]{geometry}

% ── Typography & math ──
\usepackage{fontspec}
\usepackage{amsmath,amssymb}
\usepackage{setspace}
\setstretch{1.25}

% ── Graphics ──
\usepackage{tikz}
\usetikzlibrary{arrows.meta,calc,positioning}

% ── Colors ──
\usepackage{xcolor}
\definecolor{accent}{HTML}{1a73e8}
\definecolor{darkgray}{HTML}{333333}
\definecolor{lightgray}{HTML}{f5f5f5}
\definecolor{notebg}{HTML}{e8f0fe}
\definecolor{warnbg}{HTML}{fef7e0}
\definecolor{calcbg}{HTML}{e6f4ea}
\definecolor{calcframe}{HTML}{188038}

% ── Links ──
\usepackage{hyperref}
\hypersetup{colorlinks=true,linkcolor=accent,urlcolor=accent}

% ── Layout ──
\usepackage{enumitem}
\setlist[itemize]{leftmargin=1.8em,itemsep=2pt}
\setlist[enumerate]{leftmargin=1.8em,itemsep=2pt}
\usepackage{booktabs}
\usepackage{tcolorbox}
\usepackage{titlesec}
\usepackage{fancyhdr}
\usepackage{lastpage}

% ── Header/Footer ──
\pagestyle{fancy}
\fancyhf{}
\fancyhead[L]{\small\textcolor{darkgray}{《心理学与生活》第三单元 · 发展与教育}}
\fancyhead[R]{\small\textcolor{darkgray}{六讲 · 核心概念讲义}}
\fancyfoot[C]{\small\textcolor{darkgray}{\thepage\ / \pageref{LastPage}}}
\renewcommand{\headrulewidth}{0.4pt}

% ── Section styling ──
\titleformat{\section}
  {\Large\bfseries\color{accent}}
  {\thesection}{1em}{}
\titleformat{\subsection}
  {\large\bfseries\color{darkgray}}
  {\thesubsection}{1em}{}

% ── Custom boxes ──
\newtcolorbox{keyconcept}[1][]{
  colback=notebg, colframe=accent,
  fonttitle=\bfseries, boxrule=0.5pt,
  left=8pt, right=8pt, top=6pt, bottom=6pt, title={#1}}
\newtcolorbox{exercise}[1][]{
  colback=warnbg, colframe=orange,
  fonttitle=\bfseries, boxrule=0.5pt,
  left=8pt, right=8pt, top=6pt, bottom=6pt, title={#1}}
\newtcolorbox{calcbox}[1][]{
  colback=calcbg, colframe=calcframe,
  fonttitle=\bfseries, boxrule=0.5pt,
  left=8pt, right=8pt, top=6pt, bottom=6pt, title={#1}}

\newcommand{\daytitle}[2]{%
  \section*{#1\quad #2}%
  \addcontentsline{toc}{section}{#1\quad #2}}

\begin{document}

% ═══ 封面 ═══
\begin{center}
  \vspace*{2cm}
  {\Huge\bfseries\color{accent} 发展与教育}\\[0.4cm]
  {\LARGE 六讲 · 核心概念讲义}\\[1.2cm]
  {\large 《心理学与生活》第三单元：人是怎么变成现在这个样子的}\\[0.5cm]
  {\small\textcolor{darkgray}{依据：南京大学 Coursera 课程《心理学与生活》第三单元中文字幕\\（01 延迟满足、02 婴儿气质、03 情感依恋、04 咿呀学语、\\05 性别认同、06 游戏人间、07 学习方式、08 行为塑造）}}\\[0.3cm]
  {\small\textcolor{darkgray}{\today}}
  \vfill
  {\small 共 6 讲，全书约 2.5–3 小时：先读讲解，再合上讲义把概念讲出来，最后做自测。}
\end{center}
\newpage

\tableofcontents
\newpage

% ═══ 使用说明 ═══
\section*{开始前：这本讲义在做什么}
\addcontentsline{toc}{section}{开始前：这本讲义在做什么}

第三单元回答的是两个问题：

\begin{enumerate}
  \item \textbf{人是怎么变成现在这个样子的}（第 1、2、3、5 讲）：从一颗受精卵开始，身体、大脑、气质、依恋、自我意识与性别概念，哪些是出厂自带的，哪些是环境装上去的；
  \item \textbf{学习是怎样发生的}（第 4、6 讲）：不会听课的婴儿靠什么学会一切——条件反射、习惯化、模仿与观察学习；以及奖励和惩罚这两件最常用的教育工具，边界在哪里。
\end{enumerate}

这个单元的难，不在术语多，而在于它反复纠正两类日常直觉：一类是\textbf{决定论}（“三岁看大，七岁看老”“本性难移”），另一类是\textbf{常识化归因}（“自制力就是能忍”“奖励越多动力越足”“男孩女孩天生就不一样”）。所以读法是：\textbf{先看清每个概念在反驳什么常识，再看它给出的机制是什么}。

\begin{keyconcept}[六讲要带走的三样工具]
\begin{itemize}
  \item \textbf{冷热系统旋钮}：自我控制不是硬忍，而是把注意力从诱惑（热刺激）上调开——自制力是一种可以调配的认知资源，热系统先熟、冷系统后长，所以孩子比成人冲动是发育规律，不是品行问题。
  \item \textbf{匹配观}：气质无好坏，看它与环境是否匹配；依恋源于接触安慰而非喂养；男女婴的差异很小，是被环境强化放大的——先天给原料，环境决定成品。
  \item \textbf{学习三引擎与“习得 $\neq$ 表现”}：经典条件反射、操作条件反射、观察学习；奖励与惩罚管理的是行为是否\textbf{表现}出来，而不是学没学会——这是全单元收束的主线。
\end{itemize}
\end{keyconcept}

\begin{center}
\begin{tikzpicture}[>=Stealth,node distance=0.35cm,
  must/.style={draw,rounded corners=4pt,thick,accent,fill=notebg,minimum height=0.95cm,align=center,font=\small},
  rec/.style={draw,rounded corners=4pt,thick,darkgray,fill=lightgray,minimum height=0.95cm,align=center,font=\small}]
  \node[rec] (d1) {第 1 讲\\发展心理学与延迟满足};
  \node[must,right=of d1] (d2) {第 2 讲\\身体、大脑与气质};
  \node[must,right=of d2] (d3) {第 3 讲\\依恋};
  \draw[->,thick] (d2) -- (d3);
  \node[must,below=1.2cm of d1] (d4) {第 4 讲\\条件反射、模仿与语言};
  \node[rec,right=of d4] (d5) {第 5 讲\\情绪、自我与性别};
  \node[must,right=of d5] (d6) {第 6 讲\\游戏、道德与奖惩};
  \draw[->,thick] (d4) -- (d5);
  \draw[->,thick] (d5) -- (d6);
  \draw[->,dashed,accent] (d4.south) to[out=-55,in=-125] node[midway,below,font=\scriptsize,darkgray]{“习得 $\neq$ 表现”首尾呼应} (d6.south);
\end{tikzpicture}
\end{center}

\noindent\textbf{轻重缓急}：蓝框为\textbf{必看}（第 2、3、4、6 讲是全单元的价值锚点——匹配观、学习引擎与奖惩边界都在这里），灰框为\textbf{建议}。第 1 讲独立成篇，可放在任何位置读；第 3 讲依赖第 2 讲（气质与遗传），第 5 讲依赖第 4 讲（观察学习），第 6 讲依赖第 4、5 讲（强化与游戏）。时间紧可以先走必看路径：第 2 讲 $\rightarrow$ 第 3 讲 $\rightarrow$ 第 4 讲 $\rightarrow$ 第 6 讲，其余讲次随后补；实线为内容依赖，虚线为主题呼应。

\textbf{每讲的学习流程}（全书约 2.5–3 小时）：
\begin{enumerate}
  \item 读当天的讲解，重点看“它反驳了什么常识”和“机制是什么”；
  \item 合上讲义，把当天的核心概念\textbf{讲给别人听}（或讲给自己听），卡壳的地方就是没懂的地方；
  \item 做“今日自测”，对照书末的答案。
\end{enumerate}

\textbf{边界声明}：本讲义覆盖第三单元 8 讲字幕（约 86 分钟课程视频）的全部枢纽概念，按 6 讲重新组织（原第 7 讲“学习方式”并入第 4 讲，原第 8 讲“行为塑造”并入第 6 讲）。课程中教师提及的孕育期保健、养育操作建议只保留与机制直接相关的部分；课程预告的《心理咨询与治疗》相关内容不在本讲义范围内。字幕中的旧表述（如个别现已更新的术语）在正文相应处单独注明。

\newpage

% ═══ 第 1 讲 ═══
\daytitle{第 1 讲}{人为什么会变——发展心理学与延迟满足}
\noindent\textbf{优先级}：建议\quad|\quad\textbf{难度}：$\star\star$\quad|\quad\textbf{用时}：约 20 分钟

\textbf{核心问题}：孩子能不能忍住不吃棉花糖，靠的究竟是“意志力硬扛”，还是别的什么？

\subsection*{1.1 棉花糖实验：只点一个名}

20 世纪 60 年代，斯坦福大学的沃尔特·米歇尔做了著名的\textbf{棉花糖实验}：给孩子一颗棉花糖，承诺如果等实验员回来还不吃，就再奖励一颗。多数孩子坚持不到 3 分钟，约三分之一能等上 15 分钟。近 20 年后，米歇尔跟踪了其中 600 多名已是高中生的孩子：当年能等待的孩子，在学业成绩、处理问题能力和同伴关系上普遍更好。

这个实验在本讲里只点一个名——真正值得拆的是米歇尔给出的\textbf{机制解释}。需要提醒：跟踪结果是\textbf{相关}而非因果（能等的孩子可能本来就有更好的家庭环境），这一点课程没有展开，是本讲义作者的补充。

\subsection*{1.2 冷系统与热系统：自制力是注意力的调配}

米歇尔反复观察孩子的行为，发现秘诀在于\textbf{转移注意力}：能等的孩子有的把棉花糖当成一幅画慢慢欣赏，有的对着它唱歌跳舞——他们不是忍住了不看，而是把注意力\textbf{调开}了。

\begin{keyconcept}[冷系统与热系统]
\begin{itemize}
  \item \textbf{热系统}：情绪的基础，一种先天的、“受刺激就有反应”的注意机制——它把注意力\textbf{持续固定在}诱惑物（强化物）上，推动人立刻行动。约等于日常说的“感性”。
  \item \textbf{冷系统}：认知的、策略性的注意调配，表现为各种\textbf{分心活动}——把注意力从诱惑物上移开，从而延长等待时间。约等于日常说的“理性”。
\end{itemize}
热系统发展得比冷系统\textbf{早}：人生最初几年热系统主导，所以儿童表现得冲动是发育阶段的必然，不是“不懂事”。随着年龄增长，冷系统逐渐发展起来，在延迟满足中占据优势。
\end{keyconcept}

这条机制把“自制力”从道德品质重新定义为\textbf{认知技能}：问题不是“你能不能扛住诱惑”，而是“你会不会把注意力从诱惑上调开”。想戒烟的人盯着烟灰缸硬扛，用的是热系统对热系统；起身去倒杯水、换个话题，用的才是冷系统。

\subsection*{1.3 发展心理学：变化与稳定的双重主题}

米歇尔的冷热系统会随年龄此消彼长——这正落在\textbf{发展心理学}的研究范围内：考察个体从一颗受精卵到死亡的所有生命阶段里，心理现象与规律如何展开。

\begin{keyconcept}[发展的两个面孔]
\begin{itemize}
  \item \textbf{变化}：人在不断成长、不断变化，直到生命的最后一刻；不同阶段有不同的心理品质与发展任务。
  \item \textbf{稳定}：人的行为又保持着相对稳定——所以“三岁看大，七岁看老”有观察价值，但正因为还有变化的一面，早期表现\textbf{不该被当成命运定论}。
\end{itemize}
\end{keyconcept}

这条主线会贯穿全单元：第 2、3 讲讲“哪些东西出厂就定、多稳定”，第 4、6 讲讲“哪些东西靠后天学习、怎么改变”。

\begin{keyconcept}[易错点与失效边界]
\begin{itemize}
  \item 延迟满足\textbf{不是}忍耐力竞赛，也不是道德说教——它是注意力调配的认知能力，可以用策略训练（制造分心、改变对诱惑物的表征）。
  \item 棉花糖实验的跟踪是相关研究：不能据此推出“小时候能等 $\Rightarrow$ 长大必成功”，更不能倒过来给孩子贴标签。
  \item “冷系统好、热系统坏”是误读：热系统提供情绪与动力，是生存的必需品；问题只在于它\textbf{主导}了不该由它主导的场景。
\end{itemize}
\end{keyconcept}

\begin{exercise}[今日自测]
\begin{enumerate}
  \item （辨析）“能延迟满足的孩子就是更能忍”——用冷/热系统的机制说明这句话错在哪里。
  \item （应用）孩子写作业时总忍不住看手机。依据本讲机制，给出两条不靠“没收手机硬扛”的具体策略，并说明它们各自调用了哪个系统。
  \item （讲解输出）向一位坚信“三岁看大，七岁看老”的长辈解释：这句话对在哪里、又不能推到哪里。
\end{enumerate}
\end{exercise}

\newpage

% ═══ 第 2 讲 ═══
\daytitle{第 2 讲}{天生的出厂设置——身体、大脑与气质}
\noindent\textbf{优先级}：必看\quad|\quad\textbf{难度}：$\star\star$\quad|\quad\textbf{用时}：约 25 分钟

\textbf{核心问题}：一个刚出生的婴儿，哪些是“出厂已装好”的，哪些是后来才装上的？

\subsection*{2.1 身体发展：四条规律一句话}

婴儿身体的变化有章可循，课程给出四条原则，各记一句话即可：

\begin{center}
\small
\begin{tabular}{@{}ll@{}}
\toprule
\textbf{原则} & \textbf{一句话} \\
\midrule
头尾原则 & 发展从头到脚：先会抬头转头，再会动手臂和脚。 \\
近远原则 & 发展从躯干到末端：先会动身体，再到上臂、前臂，最后手指灵活。 \\
等级整合原则 & 先掌握简单技能，再整合出复杂技能（先会拍打，再会拇指食指拿捏）。 \\
系统独立原则 & 不同系统速率不同：神经系统先快后慢，性别特征系统先慢后快。 \\
\bottomrule
\end{tabular}
\end{center}

一个直观的数字：新生儿头部占身长的四分之一，两岁时约五分之一，成年后为八分之一——“九头身”就是这么来的。

\subsection*{2.2 反射：出厂自带、限时供应的程序}

新生儿一出生就带有一批\textbf{反射}——不需要学习、由刺激自动触发的动作，多数在出生后几个月内消失。记两个有趣的例子就够：

\begin{itemize}
  \item \textbf{抓握反射}：把手指压在婴儿手掌上，他会紧紧握住。你会以为这是孩子对你的回应，从而更喜欢他——反射在无意中\textbf{促进了亲子关系}。
  \item \textbf{强直性颈反射}（“击剑反射”）：婴儿躺着时把他的头拨向一侧，同侧手脚伸出、对侧弯曲，活像一个击剑运动员。
\end{itemize}

其余反射（吸吮、吞咽、定向、游泳反射等）同理：要么服务生存，要么服务社交，都是限时供应的出厂程序。

\subsection*{2.3 大脑：数量早已封顶，发展的是连接}

\begin{keyconcept}[神经元只减不增，连接才是发展的主角]
婴儿出生时已具备人脑约 80\% 的神经元，出生后头两年基本把其余 20\% 长全——\textbf{两岁之后，神经元只会死亡、不再新生}。但不必担心：大脑的发展不靠加零件，靠\textbf{连线}。神经元之间会建起几十亿、上百亿个新连接，数量远超肌体所需；于是那些用不上的连接通过\textbf{突触修剪}被清除，常用的留下，整个网络的运作效率越来越高。
\end{keyconcept}

电影《超体》（Lucy）流行过一种说法：“人脑只开发了不到 20\%。”课程拿它当梗提起，但这其实是一个流传甚广的误读——神经连接网络始终在工作，并不存在一大片“闲置待开发”的脑区。把它当趣闻记住，别当事实引用。

\subsection*{2.4 气质：没有好坏之分的“出厂脾气”}

婴儿出生不久就显出秉性差异：课程教师举例，他的孩子丁丁可以安安静静在床上坐一个下午，摆弄一把扇子和一只橡皮鸭，乐此不疲；另一个孩子却急躁得多，扔勺子、扯纱巾、挥舞小手，一刻不停。

注意：这里的\textbf{气质}不是日常说的“某人有气质”（风度、修养），而是指\textbf{个体的唤醒模式和情绪特点}——在活动水平、易激怒水平等维度上的稳定表现，且在不同年龄阶段保持相对稳定。

\begin{keyconcept}[三种气质类型与“匹配”思想]
\begin{itemize}
  \item \textbf{易养型}：倾向积极，作息规律，适应性强，有好奇心，情绪中低强度。
  \item \textbf{难养型}：心境偏消极，适应性弱，面对新情境倾向于退缩。
  \item \textbf{发动缓慢型}：不太活跃，对环境反应平静，适应缓慢。
\end{itemize}
（许多婴儿是混合特点。）\textbf{没有哪种气质是好的或坏的}，只有哪种更能适应它所处的环境：难养型婴儿若换来父母的愤怒与责备，学龄期更可能出现行为问题；若得到温暖、关爱与包容，则大可能避免。文化也参与裁决——在某些“哭闹才喂奶”的文化里，难养型婴儿反而得到更多营养，在恶劣环境中更容易存活。
\end{keyconcept}

\subsection*{2.5 遗传：不只塑造人，还反向塑造环境}

基因检测与双生子研究表明，社会技能、领导力、敏感度、幸福感等特质受遗传影响的程度可能达 50\% 甚至更高。更值得记住的是反方向：\textbf{遗传因素会反过来决定环境}——孩子主动挑选有利于自己的环境，活泼好动的孩子被体育活动吸引，安静的孩子更可能爱上画画。教育的提示是：注意这些偏好，提供相应的培养与鼓励；同时也不必完全顺着出厂设置养，适当让孩子体验其他气质类型的活动方式，可以增强灵活性。

\begin{keyconcept}[易错点与失效边界]
\begin{itemize}
  \item “两岁后神经元只减不增”$\neq$“大脑两岁定型”：发展的是\textbf{连接}，突触修剪伴随一生，可塑性始终存在。
  \item 气质三分类是\textbf{维度上的倾向}，不是人格判决书；大量孩子是混合型。
  \item “遗传影响 50\%”说的是群体层面的\textbf{差异来源比例}，不能读成“某个人的一半由基因决定”，更不能推出“教育无用”——恰恰是基因通过环境起作用。
\end{itemize}
\end{keyconcept}

\begin{exercise}[今日自测]
\begin{enumerate}
  \item （判断）“既然神经元两岁后就只减不增，早教的意义就是趁两岁前拼命刺激大脑。”这个说法对不对？说明理由。
  \item （辨析）一位父亲说：“我家孩子天生难带，脾气坏，没办法。”用“气质 $\times$ 环境匹配”指出这句话的两个问题。
  \item （应用）你的侄女被贴了“难养型”的标签，全家都很焦虑。基于本讲内容，你会给她的父母提哪两条具体建议？
  \item （讲解输出）向朋友解释“遗传会反向塑造环境”这句话，并用“孩子挑选环境”举一个自己的例子。
\end{enumerate}
\end{exercise}

\newpage

% ═══ 第 3 讲 ═══
\daytitle{第 3 讲}{第一份关系——依恋如何塑造人}
\noindent\textbf{优先级}：必看\quad|\quad\textbf{难度}：$\star\star$\quad|\quad\textbf{用时}：约 25 分钟

\textbf{核心问题}：婴儿与养育者之间的第一份关系，凭什么影响他成年之后怎么去爱？

\subsection*{3.1 依恋：先于语言的安全感}

\textbf{依恋}是指寻求与某人的亲密，并且当其在场时感到安全的心理倾向。它不是人类专利：从习性上看，依恋是动物幼崽\textbf{适应生存的特定反应}——待在母亲身边，才有奶吃、有保护。问题是：依恋究竟靠什么建立？流行的答案（源自弗洛伊德）是“谁喂养，就依恋谁”——婴儿对母亲依恋，因为母亲提供物质满足。

\subsection*{3.2 哈洛的恒河猴：依恋不源于喂养}

20 世纪初，心理学家哈洛让刚出生的幼猴与母亲分离，转而和两个“代理母亲”生活 6 个月：一个\textbf{铁丝母亲}胸前有喂奶装置，另一个\textbf{绒布母亲}柔软舒适但没有奶。结果是：幼猴与绒布母亲待在一起的时间\textbf{多得多}，只有饿了才跑到铁丝母亲那儿喝几口，然后立刻回来紧紧抱住绒布母亲。

\begin{keyconcept}[接触性安慰]
依恋的关键因素是\textbf{身体接触}——哈洛称之为“接触性安慰”。与母亲温暖而舒适的身体接触，才是幼崽安全感的源泉；物质喂养排在其后。对养育的直接推论：不要以为让婴儿吃饱喝足就够了，想建立亲密关系、给足安全感，\textbf{多抱抱他}。
\end{keyconcept}

\subsection*{3.3 敏感期与印刻}

婴幼儿的很多发展都有\textbf{敏感期}——某项技能获得最大发展与成熟的最佳时期。洛仑兹把鸟类社会依恋的形成称为\textbf{印刻}：幼鸟在敏感期看到的第一个移动目标，会被它当成母亲并形影不离（洛仑兹就用这招让一群小鸟围着他转）。人类婴儿的依恋敏感期约从 6 个月大开始，但\textbf{何时结束尚不明确}——电影《亲爱的》里，被拐孩子在充满温暖的原生家庭长大，后来仍对收养家庭慢慢形成了新的依恋。

\subsection*{3.4 四种依恋类型：看分离与重聚}

心理学家根据儿童与父母\textbf{分离和重聚时}的表现，把依恋分为四种类型：

\begin{center}
\small
\begin{tabular}{@{}lll@{}}
\toprule
\textbf{类型} & \textbf{分离—重聚表现} & \textbf{对应的养育者互动} \\
\midrule
安全型 & 离开时轻微抗议，回来主动寻找、易被安抚 & 对孩子信号敏感，及时回应 \\
回避型（不安全） & 离开时不显痛苦，回来不主动恢复接触 & 少互动、不回应、少有身体接触 \\
抗拒型（不安全） & 离开时大哭，回来又踢打推开，表现矛盾 & 回应时有时无，不连续不稳定 \\
混乱型（不安全） & 困惑、恐惧交织，在养育者身边也害怕 & 在一定程度上忽略儿童需要 \\
\bottomrule
\end{tabular}
\end{center}

\subsection*{3.5 成人依恋：连续，但不是判决}

研究者用描述性的陈述测量成人依恋风格，结果约 60\% 的人自我报告为安全型，回避型与矛盾型各约 20\%——与婴儿依恋类型的分布\textbf{惊人地相似}。进一步的研究发现，幼年依恋质量与成年后的恋爱关系有非常大的关联。

\begin{keyconcept}[连续性 $\neq$ 决定论（争议标注）]
“成人依恋风格是否由早年依恋模式决定”这个命题\textbf{到今天仍充满争议}：依恋关系不是一次确定就不能改变，许多个体在经历成年期的人际关系（比如伴侣持续的温暖与支持）之后会改变依恋模式。可以确定的只是：无论婴儿还是成人，\textbf{情感安全}都非常重要。
\end{keyconcept}

\begin{keyconcept}[易错点与失效边界]
\begin{itemize}
  \item “管饱就行”被哈洛实验直接反驳：依恋的第一原料是\textbf{接触与回应}，不是食物。
  \item 依恋类型分布的“相似”是群体统计，不能对单个个体做“你小时候回避型，所以恋爱注定失败”的推演——决定论有争议，且成年关系可以重塑依恋。
  \item 课程开头的成人依恋自测（三组描述）只是\textbf{风格粗筛}，不是诊断工具；本讲自测已将其改造为情境判断题。
\end{itemize}
\end{keyconcept}

\begin{exercise}[今日自测]
\begin{enumerate}
  \item （辨析）用哈洛的实验设计说明：为什么“铁丝母亲有奶、绒布母亲没奶”这个对照，能反驳“依恋源于喂养”？
  \item （应用）下面三种做法，哪一种最可能帮助形成安全型依恋？为什么？（a）孩子一哭就立刻喂奶，但很少抱他；（b）对孩子的信号敏感、及时回应，多身体接触；（c）严格按时喂养和安抚，建立铁打的作息规律。
  \item （判断）“小时候形成了不安全依恋的人，成年后的亲密关系注定不幸。”判断正误并给出本讲的依据。
  \item （讲解输出）向一对新手父母解释“接触性安慰”，并把它翻译成一条今晚就能执行的育儿动作。
\end{enumerate}
\end{exercise}

\newpage
% ═══ 第 4 讲 ═══
\daytitle{第 4 讲}{学习的引擎——条件反射、模仿与语言}
\noindent\textbf{优先级}：必看（全单元核心）\quad|\quad\textbf{难度}：$\star\star\star$\quad|\quad\textbf{用时}：约 30 分钟

\textbf{核心问题}：婴儿不会听课、看不懂书，他是怎么学会这个世界的？

\subsection*{4.1 经典条件反射：中性刺激被“传染”上意义}

一个月大的迈克和家人外出遇上暴风雨，电闪雷鸣，每打一次雷他的哭声就高一级。后来，\textbf{光是闪电}就足以让他害怕得哭出来；长大后，闪电的景象仍会让他胸闷胃痛。

\begin{keyconcept}[经典条件反射]
个体学会以特定方式对一个\textbf{原本中性}的刺激做反应——因为它总和一个天然引发该反应的刺激\textbf{成对出现}。巴甫洛夫的狗是原型：铃声（中性）反复与肉（天然引起流口水）同时出现，最后狗听见铃声就流口水。
\end{keyconcept}

婴儿很早就靠它学习：每次吸吮甜水后轻敲一下脑门，很快婴儿一被敲脑门就开始吮吸；大人一边说“宝宝，吃饭”一边端出米饭，反复多次后，听到“饭”就知道要开饭了——语言习得里有大量经典条件反射的范例。

\subsection*{4.2 操作条件反射：结果被“记功”的行为会再来}

婴儿还会\textbf{为了想要的结果}故意做出行为：发现一哭父母就立刻来到身边，他便慢慢知道哭的好处，学会“工具性的哭泣”；反之，哭毫无回应，这种行为就会消退。

\begin{keyconcept}[操作条件反射]
斯金纳的老鼠：笼中老鼠无意中触碰开关得到食物，反复几次后便直奔开关——\textbf{先有反应，再有结果；行为的结果增加了该行为出现的频率}。与经典条件反射的顺序正好相反（后者是先有刺激、再有反应）。自发反应会根据结果的\textbf{正性或负性}而增强或减弱。
\end{keyconcept}

连语音都能被这样塑造：孩子要拒绝时发出“嗯——”的声调，大人便把东西拿走，这个音就“有用”了，于是被保留下来，成年后还有人用同样的声调撒娇。

\subsection*{4.3 习惯化：最安静的一种学习}

第一次看广场舞觉得新奇，天天听同一首歌就烦得想扔鸡蛋——这是\textbf{习惯化}：同一刺激重复多次呈现后，反应逐渐降低。婴儿对新刺激会先全神贯注（定向反应），呈现多次后便不再关注。习惯化让人把有限的注意力省给\textbf{新异}的东西，是过滤器的学习。

\subsection*{4.4 观察学习与“习得 $\neq$ 表现”}

班杜拉提出第三种引擎：\textbf{观察学习}。学习者不必亲自做出反应，也无需亲身体验强化，只需\textbf{观察他人（榜样）的行为及其后果}就能学会——没有人教过你怎么谈恋爱，你看多了就会了。

\begin{keyconcept}[班杜拉的玩偶实验：习得与表现分离]
儿童观看成人对人形玩偶施暴的影片，然后被带进有同款玩偶的房间：影片结尾成人\textbf{受罚}，儿童很少表现出暴力；成人\textbf{受奖}，儿童暴力行为更多；\textbf{无结果}者居中。关键在后续：隔一段较长时间再放回房间，三组儿童的暴力行为\textbf{向中间趋近}——受罚组上升，受奖组下降。
\end{keyconcept}

这说明：\textbf{仅仅通过观察，三组儿童都已经学会了暴力行为}（习得）；榜样受罚或受奖，影响的只是他们\textbf{是否表现出来}（表现）；而结果的限制作用随时间减弱，学会的东西终究会露头。这条“习得 $\neq$ 表现”的线，第 6 讲会用它来收束全单元。

\subsection*{4.5 语言：天生的机制，后天的素材}

语言学习几乎用上了上面所有引擎。模仿是主力之一：课程教师的孩子把“鸡蛋”叫成“蛋鸡”——模仿不一定照搬，会自我创造，但没有最早的“鸡蛋”，就不会有“蛋鸡”。

但“语言完全是学来的”有一个硬伤：婴幼儿在非常短的时间内就掌握了语法、语音、句法等大量信息，快得\textbf{不像逐条强化的产物}。于是语言学家乔姆斯基提出\textbf{先天论}（注意：先天论不是说婴儿天生会说话）：人生来具有学习语言的能力，所有语言共享一种\textbf{普遍语法}结构，大脑中存在一个\textbf{语言习得机制}（LAD），专门负责从周围语言里提取特征和结构。后续研究还发现了与言语产生有关的特殊基因，但尚待进一步证实。目前\textbf{没有一种观点能完美解释语言获得}，更可能的是不同因素在不同阶段各起作用。

两个常被误读的“先天”证据：
\begin{itemize}
  \item \textbf{“妈妈”全球同音}：咿呀学语期的婴儿会\textbf{自动生成}各种语言里都有的音（a、ma、ba），“ma”恰好最容易发；一开始它没有意义，只是妈妈恰好在场，用兴奋的表情、拥抱和亲吻让它“有了意义”。所以趋同源于\textbf{婴儿的发音生理}，而非什么神秘的统一词根。
  \item \textbf{6 个月的分水岭}：6 个月后，婴儿开始难以分辨母语中不用的发音——他被自己的母语环境“稳定”下来。窗口的存在支持先天机制，窗口关闭的方向由环境决定：先天给机器，环境给素材。
\end{itemize}

补充两笔：聋哑儿会自发用手势表达，且手势激活的脑区与言语生成几乎相同——口语可能是从手语进化来的；大人对婴儿说话时那种音调变高、语调夸张、词语重复的\textbf{婴儿指向言语}，对语言学习极其重要——跟孩子说话，蹲下来，用他的语言。

\begin{keyconcept}[易错点与失效边界]
\begin{itemize}
  \item 两种条件反射的\textbf{顺序}别记反：经典是“刺激 $\rightarrow$ 反应”（铃声引起流口水），操作是“反应 $\rightarrow$ 结果”（按压开关带来食物）。
  \item “奖励多的行为就是学会的、没表现出的就是没学会”——班杜拉的趋同结果恰恰否定这一点：惩罚压住的只是\textbf{表现}。
  \item 先天论 $\neq$ 语言天注定：没有语言环境，LAD 无米下锅；学习论与先天论之争至今没有终审判决。
\end{itemize}
\end{keyconcept}

\begin{exercise}[今日自测]
\begin{enumerate}
  \item （辨析）下面各是哪种学习引擎？（a）听到上课铃就心跳加快；（b）为了得到小红花而举手发言；（c）看表哥骑自行车，第一次上车就会蹬；（d）对楼下的装修声渐渐充耳不闻。
  \item （应用）一个一岁半的孩子学会用尖叫让大人妥协。用操作条件反射解释这个行为是怎么被“养大”的，并给出一条打断它的做法。
  \item （判断）“只要在孩子面前惩罚了施暴的榜样，孩子就不会学会暴力行为。”结合班杜拉实验的后续结果判断。
  \item （讲解输出）向朋友解释“语言习得机制”假说要解释的难题是什么，以及“妈妈同音”为什么不能算它的证据。
\end{enumerate}
\end{exercise}

\newpage

% ═══ 第 5 讲 ═══
\daytitle{第 5 讲}{情绪、自我与性别——社会化的孩子}
\noindent\textbf{优先级}：建议\quad|\quad\textbf{难度}：$\star\star$\quad|\quad\textbf{用时}：约 25 分钟

\textbf{核心问题}：孩子的情绪、“我”的概念和“男孩/女孩”的概念，分别是从哪里来的？

\subsection*{5.1 情绪：天生的表情，后天的扩充}

婴儿有与成人相差无几的表情——微笑、生气、皱眉，且基本表情在不同文化间\textbf{基本相似}。发展心理学家伊扎德认为，婴儿天生有一套反映基本情绪状态的表情；随着主导情绪的前额叶与边缘系统成熟，这套表情被不断扩展、修正，情绪体验越来越广。约 1 至 2 岁时出现\textbf{社会性微笑}：对母亲等照料者的微笑比对物体更频繁，而且\textbf{得不到回应就会减少}——注意，这正是第 4 讲操作条件反射在情绪上的运作：微笑因为“有用”而被保留。

\subsection*{5.2 社会性参照：看别人脸色定自己反应}

\begin{keyconcept}[社会性参照]
约 8、9 个月起，婴儿在\textbf{不确定}的情境中，会有意识地搜寻他人的情感信息来决定自己如何反应——这就是社会性参照。再大一点，来自不同人的\textbf{相互矛盾}的信息会使他不安：孩子用勺子敲桌子，妈妈生气、爸爸却笑，孩子收到的信号是混乱的——他多半会继续敲，来测试到底行不行。推论很直接：想教育孩子，\textbf{养育者的态度应当一致}。
\end{keyconcept}

\subsection*{5.3 自我意识：镜子里的那个“我”}

我们不是天生知道“镜中人是我”。研究者用一个极简实验检验：趁婴儿不注意，在他鼻子上点一个红点，再抱到镜子前——如果他伸手去摸\textbf{自己的}鼻子，说明他认出了自己，已有初步的自我意识；如果他摸的是\textbf{镜中}那个人的鼻子，说明还没有。这就是著名的\textbf{点红实验}：自我意识大约在 12 个月左右开始发展（字幕未给出该实验的研究者姓名，本讲义不杜撰归属）。

\subsection*{5.4 sex 与 gender：先把两个词分开}

英文表格里“性别”一栏常写 gender 而不是 sex，因为两者不同：\textbf{sex} 指生理、解剖学上的性别特征；\textbf{gender} 指两性被社会期待扮演的角色，以及受文化影响较大的那些东西。分开之后才能看清下一个问题：性别差异里，哪部分是 sex，哪部分是 gender？

\subsection*{5.5 性别差异：起点很小，被环境放大}

\begin{keyconcept}[差异主要被环境强化放大]
研究发现，即便存在生理差异，\textbf{女婴和男婴之间的差异其实很小}——光看外表，成人很难区分婴儿的性别。男孩拿到玩具枪喊“冲啊”，女孩抱布娃娃说“我给你看病”，本质上\textbf{都是角色扮演的模仿游戏}，玩具之间没有高下，是我们人为做了定义。差异会\textbf{随年龄不断放大}，而放大器是环境：男孩玩布娃娃受到的阻力，远大于女孩玩玩具枪；男孩被鼓励探索世界，女孩得到更多拥抱、被留在更近的距离。
\end{keyconcept}

学前期的孩子甚至会\textbf{比成人更刻板}：小男孩穿粉色衬衫可能被全班嘲笑。此期孩子也开始偏好同性伙伴——而这偏好本身，可能正是“我们过分强调男女差异”的产物，因为人总爱和相似的人玩。

怎么解释性别示意行为的形成？\textbf{激素理论}（性激素）解释不全面；得到普遍支持的是\textbf{社会学习理论}：孩子通过观察父母、老师、兄弟姐妹、玩伴来学习性别行为，肯定与鼓励不断强化，书籍与媒体持续引导（这正是第 4 讲观察学习的应用）。课程还提到，若此过程中缺乏可模仿的对象、未能建立恰当的联系，可能出现课程原文所称的“易性癖、变性人”等情况——这是课程的旧时表述，今天的规范说法是\textbf{性别认同与出生时被指定的性别不一致}（跨性别），且它是否如此形成并无定论，本讲义仅按课程口径转述。最新的教育理念倾向于\textbf{弱化性别教育}：不再强调“男孩该怎样、女孩该怎样”，而是让男性和女性都成为刚柔并济的独立个体。

\begin{keyconcept}[易错点与失效边界]
\begin{itemize}
  \item 别拿成人的性别差异去倒推婴儿：成人的差异里叠了一二十年的环境强化，不能当作“天生如此”的证据。
  \item 社会学习理论“得到普遍支持”不等于激素不起作用——课程明确说激素理论不能\textbf{完全}解释，两者是互补而非互斥。
  \item 点红实验测的是\textbf{视觉自我识别}，通过它不等于拥有完整的自我概念；未通过也不等于“没有自我”。
\end{itemize}
\end{keyconcept}

\begin{exercise}[今日自测]
\begin{enumerate}
  \item （辨析）“男孩天生好动、女孩天生文静，差异是生理决定的。”用本讲的证据指出这个推论在哪一步跨过了数据的边界。
  \item （应用）幼儿园里，男孩抢了别人的玩具，老师说“男孩子就是皮”；女孩不敢举手发言，老师说“女孩子就是害羞”。从社会学习理论看，这两句“解释”本身在起什么作用？
  \item （判断）“婴儿没有通过点红实验，说明他还没有任何自我意识。”判断并说明边界。
  \item （讲解输出）向朋友解释 sex 与 gender 的区别，并用“玩具枪与布娃娃”的例子说明“差异被环境放大”。
\end{enumerate}
\end{exercise}

\newpage

% ═══ 第 6 讲 ═══
\daytitle{第 6 讲}{游戏、道德与奖惩的边界}
\noindent\textbf{优先级}：必看\quad|\quad\textbf{难度}：$\star\star\star$\quad|\quad\textbf{用时}：约 35 分钟

\textbf{核心问题}：游戏、道德和奖惩——教育工具箱里最常用的三件工具，各自的边界在哪里？

\subsection*{6.1 自言自语与幻想游戏：思维的脚手架}

很多小孩喜欢和自己说话，别人听不懂，自己说得不亦乐乎。维果斯基认为这种\textbf{自言自语}会被用来指导儿童的行为和思维、产生新想法：成人思考难题时心里那个“内部对话”，正是以童年的自言自语为先兆；对自己说“这样不好、不好”来控制愤怒，用的也是它。

游戏则沿着一条升级线展开：
\begin{itemize}
  \item \textbf{感觉运动游戏}：简单重复的动作，看到妈妈洗盘子，自己也把手伸进水里。
  \item \textbf{符号游戏}：能通过回忆再现行为，把一种物体\textbf{假想为}别的物体——课程教师发现孩子反复做一个动作，后来才明白那是在模仿打开翻盖手机。没有道具时，儿童能发明道具或赋予日常物品新用途（垃圾桶、锅碗瓢盆皆可成玩具），所以教育专家主张给孩子\textbf{灵活、限制少}的游戏材料。
  \item \textbf{幻想游戏}：以想象为基础，让物体和人物获得新身份、新意义——过家家是典型。它不只是好玩：儿童借此理解社会与现实、扮演爸爸妈妈医生超人，\textbf{抽象能力}在游戏中被锻炼。
\end{itemize}

\subsection*{6.2 游戏与道德：在冲突里学会做人}

随着年龄增长，游戏中的互动与冲突增加，而\textbf{冲突及其解决}为道德意识与行为提供了重要条件——弗洛伊德认为游戏是儿童形成道德的最好方式之一。道德不是听课听来的，是和玩伴抢玩具、和好、再分工的过程里长出来的。

\subsection*{6.3 科尔伯格：道德推理的三水平}

课程给出经典两难：海因茨的妻子患癌症，唯一救命的药被药商抬到高价，海因茨凑不够钱，夜里潜入药店偷了药。\textbf{他该不该？}心理学家科尔伯格据此把道德发展分为三个水平——注意，定水平的依据是\textbf{推理的理由}，不是“偷/不偷”的结论：

\begin{center}
\small
\begin{tabular}{@{}ll@{}}
\toprule
\textbf{水平} & \textbf{典型推理（同一个结论，不同的依据）} \\
\midrule
前习俗水平 & “不偷妻子会死”或“偷了会坐牢”——看结果与惩罚； \\
            & 或“药商就是赚钱，活该”——看功用与互惠及利。 \\
习俗水平 & “动机虽好但违法，不该偷”——看“好孩子”期待与人际和谐； \\
            & 或“合情但违法，应禁止”——看规则与法律，情理法常打架。 \\
后习俗水平 & “法律只是社会契约，可以因多数人的要求改变，动机好可减刑” \\
            & ——看契约与道义责任；更高者把生命置于普遍伦理之上。 \\
\bottomrule
\end{tabular}
\end{center}

两个边界：道德不靠一个故事形成；且海因茨测的只是\textbf{道德认知}，完整的道德还包括\textbf{道德情感与道德行为}——所以要让孩子多参与真实的人际互动与冲突，在情绪体验中形成协调与化解冲突的能力。

\subsection*{6.4 表扬与鼓励：一句话之差，方向相反}

心理学家 Dweck 的研究区分了两种口头奖励：孩子完成同样的任务后，\textbf{表扬}组听到“你一定很聪明”，\textbf{鼓励}组听到“你一定很努力”。后续分化：第二轮，被表扬的孩子只愿选与原题同难度的题——他要\textbf{保住“聪明”}；90\% 被鼓励的孩子选了更难的题。第三轮任务很难，鼓励组乐在其中、注意力集中，表扬组相当痛苦——他们开始怀疑自己是不是真聪明。第四轮难度与第一轮相同：\textbf{表扬组成绩下降 20\%，鼓励组提高 30\%}。

\begin{keyconcept}[夸努力，不夸聪明]
鼓励让孩子发现成功与否\textbf{可以由自己控制}（努力是自己的变量）；简单的表扬则剥夺了这种可控性——“聪明”是固定标签，一旦失败就是标签的崩塌。奖励不是只要说“你真聪明、你好棒”就可以的。
\end{keyconcept}

\subsection*{6.5 外在奖励会挤出内在动机}

惩罚不只有批评和体罚：\textbf{取消一个人已经习惯拥有的东西}也是惩罚——所以从孩子手里拿走玩具才那么痛苦。奖励也不只有物质和表扬，而物质奖励有时反而坏事。课程讲了一个流传的故事（\textbf{轶事级证据}，当寓言用，别当实验数据）：意大利移民萨尔瓦多爱在院子里大放歌剧，邻居们厌恶歌剧，天天来揍他，警告无效。老人改了策略：每次挨完打，给每人 10 美元；一周后降到 1 美元，邻居们愤怒了——“一美元能干什么？”从此再没人来打他，歌剧声自由了。

\begin{keyconcept}[挤出效应]
外在的物质奖励有时会\textbf{替换内在动机}：邻居们打人本来是为了“让他别放歌剧”，钱把动机换成了“为美元而打”，美元一缩水，行为便瓦解。推论：别以为加薪是唯一的激励——它可能让员工失去自发的热情。
\end{keyconcept}

\subsection*{6.6 收束：习得 $\neq$ 表现}

现在把全单元焊起来。第 4 讲班杜拉的实验已经证明：奖惩管理的是\textbf{表现}，学习早已通过观察完成。本讲的两块证据是同一枚硬币的两面：表扬（一种奖励）改变了孩子\textbf{表现出}挑战意愿的方式；外在奖励把行为的\textbf{动机账户}从内在挪到外在。所以奖惩的真实身份是“表现的调节器”，而不是“学习的发动机”——用它之前先问：我想改变的是孩子会不会，还是他愿不愿表现出来？前者靠第 4 讲的三引擎与示范，后者才轮到奖惩，而且措辞（鼓励而非表扬）、形式（内在动机的保护）都有讲究。

最后借用课程的结语：发展与教育的所有知识\textbf{都不是绝对真理}，心理学家养孩子也充满矛盾，只是他们更灵活地运用知识——遇到问题不照搬理论，结合实际去调整修正，\textbf{灵活性与积极性才是最重要的}。

\begin{keyconcept}[易错点与失效边界]
\begin{itemize}
  \item 科尔伯格的水平按\textbf{推理依据}划分，不按结论：说“该偷”可能是前习俗（怕妻子死），也可能是后习俗（生命高于财产权）。
  \item “夸孩子聪明总没错”被数据直接反驳：表扬组第四轮成绩下降 20\%。要夸就夸\textbf{可控的东西}——努力、策略、坚持。
  \item 萨尔瓦多故事是轶事：它漂亮地演示了挤出效应，但不能替代实验证据；挤出的发生有边界（本来就毫无内在动机的行为，外在奖励仍是有效的）。
  \item 惩罚的形式被低估了：拿走已习惯的东西就是惩罚，使用时要知道自己在用惩罚。
\end{itemize}
\end{keyconcept}

\begin{exercise}[今日自测]
\begin{enumerate}
  \item （辨析）同样是“海因茨不该偷药”，给出分别属于前习俗、习俗、后习俗水平的三种推理，并说明定水平的依据是什么。
  \item （应用）孩子数学考好了，你要说一句话奖励他。写出表扬版和鼓励版两句话，并预测半年内它们各自可能塑造出什么行为模式。
  \item （判断）“孩子本来不爱练琴，那就更该用物质奖励把兴趣培养起来。”用挤出效应判断并说明边界。
  \item （讲解输出/总检验）合上讲义，向朋友讲清本单元主线：从冷热系统、气质 $\times$ 环境、依恋、学习三引擎，到“习得 $\neq$ 表现”——并画出这六讲的知识地图。
\end{enumerate}
\end{exercise}

\newpage

% ═══ 参考答案 ═══
\section*{参考答案}
\addcontentsline{toc}{section}{参考答案}

\textbf{第 1 讲}
\begin{enumerate}
  \item “能忍”把自制力放在热系统对热系统的消耗战上；米歇尔的观察表明，能等的孩子靠的是\textbf{冷系统}的策略性注意调配——把棉花糖当画看、唱歌跳舞，把注意力从强化物上移开。错在把认知技能当成了意志力品质。
  \item 要点：调用冷系统制造\textbf{分心或重构}——把手机放到另一个房间（削弱热系统的刺激输入）、把作业拆成有即时反馈的小段（给冷系统可调配的目标）、约定“写完这一段看 5 分钟”（用规则替代现场搏斗）。凡是“坐在手机旁边硬扛”的方案，都是让冷系统单挑热系统，长期必败。
  \item 对的地方：发展有\textbf{稳定性}，早期表现与日后表现相关，值得观察。不能推到的地方：发展同时是\textbf{持续变化}的，冷热系统此消彼长、环境与教育持续作用，早期表现不是命运定论——把观察价值升级成判决，就跨过了证据的边界。
\end{enumerate}

\textbf{第 2 讲}
\begin{enumerate}
  \item 不对。两岁封顶的是\textbf{神经元数量}，发展的是\textbf{连接}：连接终生在建立、被突触修剪优化，可塑性不随两岁关门。早教有意义，但意义不是“抢在 deadline 前堆刺激”。
  \item 两个问题：（1）气质无好坏，“难带”不等于“脾气坏”，难养型只是适应性弱、易退缩的\textbf{倾向}；（2）结果取决于\textbf{匹配}——同样的难养型，遇到愤怒责备更可能出问题，遇到温暖包容则大可能避免。“没办法”恰恰放弃了父母手里最大的变量。
  \item 要点：（1）识别并尊重孩子的节奏，对新情境给足缓冲，不以愤怒回应消极反应；（2）提供温暖、关爱与包容的稳定互动，同时适当安排多元体验以增强灵活性。核心是把“这孩子有问题”改写成“这孩子的气质需要这样的环境”。
  \item 要点：遗传不只决定人，还通过人的\textbf{选择}决定环境——活泼的孩子被体育活动吸引、安静的孩子爱上画画，环境再反过来强化特质，形成回路。例子要求具体：某次自己主动选择的圈子/活动如何塑造了后来的自己。
\end{enumerate}

\textbf{第 3 讲}
\begin{enumerate}
  \item 对照的精妙在于\textbf{分离变量}：铁丝母亲提供食物但不提供柔软接触，绒布母亲相反。若依恋源于喂养，幼猴应更多待在铁丝母亲处；事实是它们只去喝奶，其余时间紧抱绒布母亲——喂养假设被证伪，接触性安慰胜出。
  \item （b）。安全型依恋对应的是养育者\textbf{对信号敏感、及时回应、有身体接触}。（a）只有喂养没有接触与回应，落入哈洛反驳的假设；（c）规律本身不坏，但“铁打的作息”若意味着对信号的延迟回应，更接近不安全型依恋的养育风格。
  \item 错误。依据两条：（1）“早年依恋是否决定成人依恋”至今\textbf{有争议}；（2）依恋模式可在成年期的人际关系中被重塑——持续温暖支持的伴侣关系能改变不安全依恋。幼年质量与成年恋爱是\textbf{关联}，不是判决。
  \item 要点：接触性安慰 = 温暖舒适的身体接触是安全感的源泉，优先级在喂养之上。今晚的动作示例：睡前 10 分钟肌肤相亲的拥抱与安抚，且对孩子的哭声给出\textbf{回应}而非任其“哭累”。
\end{enumerate}

\textbf{第 4 讲}
\begin{enumerate}
  \item （a）经典条件反射（铃声与“该紧张的情境”反复配对）；（b）操作条件反射（发言的结果——小红花——增加发言频率）；（c）观察学习（未亲身体验，看榜样即习得）；（d）习惯化（重复呈现，反应降低）。
  \item 养成路径：尖叫 $\rightarrow$ 大人妥协（正性结果）$\rightarrow$ 尖叫频率上升——标准的操作条件反射。打断方式：让尖叫\textbf{不再带来结果}（消退），同时给替代表达（用词语提出请求）以正性结果——注意只对行为做，不把情绪本身当错。
  \item 错误。惩罚榜样压住的只是\textbf{表现}：三组儿童都通过观察学会了暴力，随时间推移受罚组的暴力水平回升、三组趋近。想真正降低学会的内容，要管理孩子\textbf{看到什么榜样}，而不只是榜样的下场。
  \item 要点：学习论无法解释婴幼儿在极短时间内掌握语法、语音、句法的“超越式”速度，LAD 假说就是为此提出的（普遍语法 + 专门机制）。“妈妈”同音的解释在\textbf{发音生理}：婴儿咿呀期自动发出 ma 这类最易发的音，语义是妈妈在场赋予的——它不需要、也不支持普遍语法。
\end{enumerate}

\textbf{第 5 讲}
\begin{enumerate}
  \item 数据只到“婴儿期差异很小、成人难以凭外表分辨”为止；成人的明显差异里叠了一二十年的环境强化（玩具、鼓励方向、距离管理、同伴压力）。从成人差异直接倒推“生理决定”，跨过的正是\textbf{环境放大}这一步。
  \item 这两句话本身就是\textbf{社会学习理论里的强化与榜样信号}：它们把刻板期待喂给孩子，参与\textbf{制造}而不是描述差异——男孩学到“皮是被允许的”，女孩学到“沉默是被期待的”。
  \item 不严谨。点红实验测的是\textbf{视觉自我识别}这一条通道；未通过只能说明镜中识别尚未出现，不能推出“没有任何自我意识”。它只是自我意识发展的一个\textbf{指标}，不是全量定义。
  \item 要点：sex 是生理性别，gender 是社会期待的角色。玩具枪与布娃娃本质都是角色扮演的模仿，差异不在玩具而在我们给玩具贴的性别标签及后续的差别对待——差异的种子小，环境的放大器大。
\end{enumerate}

\textbf{第 6 讲（总检验）}
\begin{enumerate}
  \item 示例：前习俗——“不该偷，会被抓去坐牢”（惩罚取向）；习俗——“偷窃违法，合法公民不该这么做”（规则与秩序取向）；后习俗——“不该偷，因为应推动改变允许药商漫天要价的规则，而不是用违法对抗”（契约与普遍原则取向）。依据是\textbf{推理的理由}：同一结论可落在三个水平。
  \item 表扬版：“你真聪明。”鼓励版：“你这次很用心，难题也没放弃。”预测：前者把“聪明”变成需要保护的固定标签，孩子倾向回避挑战以保住标签；后者把成功归因于\textbf{可控变量}（努力、策略），孩子更敢选难题。Dweck 的数据：第四轮同难度下，表扬组降 20\%、鼓励组升 30\%。
  \item 基本不成立。挤出效应表明：外在奖励可能替换内在动机，奖励一撤行为即瓦解——它适合“本来就没有动机”的机械任务，不适合“想培养兴趣”的场景。培养兴趣的正道是保护自主性（选择、胜任感、意义感），奖励若用，宜轻、宜意外、宜指向努力而非交易。
  \item 地图参考：冷热系统（自制力 = 注意调配）$\rightarrow$ 出厂设置（发展四原则、反射、突触修剪、气质 $\times$ 环境、遗传选环境）$\rightarrow$ 依恋（接触性安慰、敏感期、四类型、连续性有争议）$\rightarrow$ 学习三引擎（经典/操作/观察 + 习惯化）+ 语言之争 $\rightarrow$ 社会化（参照、点红、sex vs gender、社会学习）$\rightarrow$ 游戏与道德（维果斯基、科尔伯格）+ 奖惩边界（表扬 vs 鼓励、挤出效应），全单元以“习得 $\neq$ 表现”收束：奖惩管表现，学习靠引擎。
\end{enumerate}

\vspace{1em}
\begin{keyconcept}[学完六讲后，你应该能讲给别人听的六句话]
\begin{enumerate}
  \item 自制力不是硬忍，是把注意力从诱惑上调开——冷系统管策略，热系统管冲动，孩子冲动是因为冷系统还没长好。
  \item 气质没有好坏，只有与环境的匹配；神经元两岁后只减不增，但大脑的发展靠的是连接，不是数量。
  \item 依恋源于接触与回应，不源于喂养——多抱抱，比管饱更要紧；早年依恋与成年关系有关联，但决定论有争议。
  \item 婴儿靠三台引擎学习：经典条件反射、操作条件反射、观察学习；惩罚压得住表现，压不住习得。
  \item 男女婴的差异很小，是被玩具、期待和榜样放大的；sex 是生理，gender 是剧本。
  \item 夸努力不夸聪明，外在奖励会挤出内在动机；奖惩是“表现的调节器”，不是“学习的发动机”。
\end{enumerate}
\end{keyconcept}

\begin{calcbox}[后续学习指引（本讲义未展开的部分）]
\begin{itemize}
  \item 课程第一、二单元（感知觉、记忆、思维等基础主题）与第四单元之后的章节，可按同样方式阅读。
  \item 课程预告的《心理咨询与治疗》：依恋模式改变、行为塑造的临床应用将在那里展开。
  \item 想深挖的读者：米歇尔关于冷热系统的著作、班杜拉社会学习理论、科尔伯格道德发展阶段、Dweck 的成长型思维研究，都值得回到原始文献。
\end{itemize}
\end{calcbox}

\end{document}
